% 模型假设
\section{模型假设}
\label{sec:assume}

为使问题可解并保证模型与题目规则一致，作如下假设。

\begin{enumerate}
\item \textbf{飞行运动学假设}：无人机垂直起飞与降落，爬升与下降速度恒为
$v_{\mathrm{up}}$ 与 $v_{\mathrm{down}}$；水平巡航按航线匀速飞行，速度为
$v_c$。航线巡航高度取航段上各点地面高程的最大值再加 50\,m 安全余量，即
$h_{\mathrm{cruise}}=\max_{x\in leg}\mathrm{DEM}(x)+50$\,m；服务区内作业高度
取该区地面高程，按悬停 30\,s 投递计入时延。飞行时间按三段相加，即
$t=h^+/v_{\mathrm{up}}+d/v_c+h^-/v_{\mathrm{down}}$。

\item \textbf{载荷能耗假设}：水平飞行单位里程能耗由等效航程函数 $L_g(q)$
刻画，该函数随载荷 $q$ 增大而缩短；水平能耗等于可用电量乘以航程与等效航程
之比，即 $E_{\mathrm{hor}}=E_{\mathrm{use}}\cdot d/L_g(q)$。爬升段为克服重力
做功，即 $E_{\mathrm{up}}=(m_0+q)gh^+/(\eta\cdot3.6\times10^6)$\,kWh，下降段
近似忽略。任意架次总能耗不得超过返航安全上限，即
$E_{\mathrm{total}}\le(1-\rho)E_{\mathrm{use}}$，其中 $\rho=20\%$ 为返航安全
余量。

\item \textbf{电池与充电假设}：每个架次出发前在 O01 更换满电电池；返航后电池
按两阶段曲线充电，即荷电状态低于 0.9 时充电时长为
$T\cdot(0.65\cdot(0.9-SOC)/0.9+0.35)$，否则为
$T\cdot0.35\cdot(1-SOC)/0.1$；充电完成的电池才可再次出勤，同一块电池同一时刻
仅供一架无人机使用。

\item \textbf{通信传播假设}：信号按自由空间传播衰减，附加系统损耗
$L_{\mathrm{sys}}=3$\,dB 与地形遮挡损耗 $L_{\mathrm{obs}}=10$\,dB；接收灵敏度
$P_{\mathrm{sens}}=-98$\,dBm，解调余量 $M=8$\,dB，即接收门限
$P_{\mathrm{th}}=-90$\,dBm。判定通信有效需链路总损耗不超过直连、接入、回程
三档上限（122、116、126\,dB）之一，且收发两端间不存在地形遮挡，即在连线
上以 30\,m 步长采样，任一点地形高于连线则判遮挡。

\item \textbf{中继假设}：中继无人机由 O01 起飞，悬停于候选位置执行转发任务，
悬停高度按与地面链路衔接要求取布设点高程以上一定高度且不超过 300\,m；每次
任务含 180\,s 起飞准备、30\,s 建链、服务时段与转场时间；中继单架次能源预算为
$(1-\rho)E_{\mathrm{use}}=2.56$\,kWh，超出则须返场充电再出勤。

\item \textbf{资源约束假设}：每架无人机任一时刻只能执行一个任务；货箱不可
分包，即一个货箱由单一架次一次送达；暂不考虑气象、禁飞区与无人机故障，
故障冗余问题见第\ref{sec:p4}节。
\end{enumerate}